\problemname{Trafikverket's plan}
During winter, the amount of road accidents in Lönnköping is far too great due to ice roads. 
Since the municipality of Lönnköping wants to appear even more efficient than the municipality of
Lund (see \textit{Illuminated City}), their road network consists of $N-1$ between the $N$ houses of the city.
The network is constructed such that you can drive between each pair of houses using the roads.
In other words, the road network is a tree. The Swedish Authority of transportation (Trafikverket)
has devised a brilliant plan to reduce the number of accidents: they plan to make each road in the
road network one-way, while still ensuring that everyone can get to work in the morning. In the
evening, they will reverse all directions so that everyone can get home.

The only problem is that they do not know whether their plan will work. What if some residents
can't get to work in the morning? Trafikverket have collected a list of where each resident lives
and works and want to know for each one of the them whether they can get to work given their new
plan. Their will include a description of the current road network and which direction they
plan to direct each road.

\section*{Input}
The first line of the input contains the integer $N$ ($1 \le N \le 3 \cdot 10^5$), the number of
houses in Lönnköping.

The following $N-1$ lines each contain the integers $a,b$ ($1 \leq a,b \leq N$, $a \neq b$),
meaning that the road between house $a$ and $b$ will be directed from $a$ to $b$. It is
guaranteed that the the undirected edges form a tree.

The next line contains the integer $Q$ ($1 \leq Q \leq 10^5$), the number of residents of Lönnköping.

Afterwards, $Q$ lines following, each containing the integers $H_i,W_i$ ($1 \leq H_i,W_i \leq N$),
the index of the house where resident $i$ lives and works respectively.

\section*{Output}
For each resident, print ``ja'' if they can get to work under the new plan. Otherwise,
print ``nej''. Print these in the same order of people as in the input.

\section*{Points}
Your solution will be tested on several test case groups.
To get the points for a group, it must pass all the test cases in the group.

\noindent
\begin{tabular}{| l | l | p{12cm} |}
  \hline
  \textbf{Group} & \textbf{Point value} & \textbf{Constraints} \\ \hline
  $1$    & $10$       & $N, Q \leq 1000$ \\ \hline
  $2$    & $10$       & $W_i=1$ for all $i$. \\ \hline
  $3$    & $10$       & $|a-b|=1$ for all pairs of $(a,b)$, that is to say that the tree is a line. \\ \hline
  $4$    & $25$       & För all pairs $(a,b)$ it holds that $\lfloor\frac{max(a,b)}{2} \rfloor=min(a,b)$.
  In other words, the tree is a full binary tree. \\ \hline
  $5$    & $15$       & $N \leq 50\,000$ \\ \hline
  $6$    & $30$       & No additional constraints. \\ \hline
\end{tabular}

\section*{Samples}
Sample 3 satisfies the constraints for subtask 3 and sample 4 satisfies the constraints for subtask 4.
